\chapter*{Chapter 5}
\markboth{Chapter 5}{5 Testing, Deployment and Conclusion}
\addcontentsline{toc}{chapter}{5 Testing, Deployment and Conclusion}

\textsc{\textbf{\Huge Testing, Deployment and Conclusion}} \\ \\

\stepcounter{chapter}

% ============================================================
\section*{Introduction}
\addcontentsline{toc}{section}{Introduction}

This chapter covers the final phase of the project, addressing
Sprints 7 and 8 which focused on testing, deployment, and
concluding the development work. While the previous chapters
described the implementation of features, this chapter explains
how those features were validated and prepared for production
use. The testing phase ensured that the platform functioned
correctly and met quality standards. The deployment phase
made the platform accessible to users. The conclusion summarizes
the entire project and identifies areas for future enhancement.

The work in these final sprints was different from the earlier
implementation sprints. Rather than building new features, the
focus shifted to verifying existing functionality, fixing issues
discovered during testing, and preparing the application for
production deployment. This phase required different skills and
approaches compared to feature development.

% ============================================================
\section*{Sprint 7: Testing and Quality Assurance}
\addcontentsline{toc}{section}{Sprint 7: Testing and Quality Assurance}

\subsection*{Sprint Overview}

Sprint 7 concentrated on verifying that the platform worked
correctly and met quality standards. Testing is a critical
phase in any software project that ensures features function
as intended and that defects are identified and resolved before
users encounter them. This sprint implemented a comprehensive
testing strategy covering multiple levels of verification.

The testing approach followed the principle of the testing
pyramid, with many unit tests at the base for individual
components, fewer integration tests for verifying component
interactions, and fewer still end-to-end tests for critical
user workflows. This distribution ensures that most testing
is fast and focused while still validating overall system
behavior.

The Laravel framework provided excellent support for testing
through its built-in testing tools. PHPUnit served as the
testing framework, with Laravel's HTTP tests for API endpoints,
feature tests for business logic, and unit tests for individual
methods. The test database configuration allowed tests to run
against an in-memory database that was fresh for each test.

Unit tests targeted individual methods and functions to verify
they produced correct outputs for given inputs. These tests were
fast to execute and easy to write, making them ideal for
catching bugs early in development. The authentication module
had extensive unit tests covering password hashing, token
generation, and role checking.

Feature tests verified larger pieces of functionality by
simulating user interactions. These tests made HTTP requests
to the API and verified the responses, checking both the
status codes and the content of returned data. Feature tests
for the product search ensured that queries returned expected
results with proper filtering and sorting.

The code coverage tools measured how much of the codebase
was exercised by tests. While complete coverage is not
achievable or necessary, targeting important business logic
ensured that critical paths were validated. The authentication,
search, and pricing calculations received particular attention
due to their central role in the platform.

\begin{table}[H]
\begin{center}
\begin{tabular}{|m{3.5cm}|m{5.5cm}|}
\hline
\textbf{Test Type} & \textbf{Purpose} \\ \hline
Unit Tests & Verify individual methods and functions \\ \hline
Feature Tests & Test API endpoints and business workflows \\ \hline
Integration Tests & Verify component interactions \\ \hline
Regression Tests & Ensure fixes do not break existing features \\ \hline
\end{tabular}
\caption{Testing Types Used in Sprint 7}
\label{tab:testing_types}
\end{center}
\end{table}

The test-driven development approach was used for some new
features where tests were written before the implementation.
This practice helped clarify requirements and ensured that
the implementation met the defined expectations. For existing
features, tests were written to capture current behavior and
then expanded as issues were discovered.

Bug tracking integrated with the sprint workflow. Issues
identified during testing were documented, prioritized, and
resolved before proceeding to deployment. Critical bugs that
affected core functionality received immediate attention,
while minor issues were scheduled for resolution in subsequent
iterations.

Security testing verified that the authentication and authorization
systems functioned correctly. Password requirements were validated,
session management was tested, and role-based access controls
were verified for each user type. This testing ensured that
users could only access features appropriate to their role.

Performance testing evaluated how the system handled load.
Search queries were tested with large datasets to ensure
response times remained acceptable. The database was checked
for proper indexing, and slow queries were optimized. Caching
strategies were verified to confirm they provided the expected
performance improvements.

Testing results were positive overall. The authentication
system passed all security tests. Search functionality returned
accurate results within acceptable timeframes. The role-based
access controls worked correctly for all user types. Minor
issues discovered were documented and resolved.

\begin{figure}[H]
\begin{center}
\fbox{Add screenshot: Test results / Postman API tests / Coverage report}
\end{center}
caption{Testing Results}
\label{fig:testing_results}
\end{figure}

% ============================================================
\section*{Sprint 8: Deployment and Finalization}
\addcontentsline{toc}{section}{Sprint 8: Deployment and Finalization}

\subsection*{Sprint Overview}

Sprint 8 focused on preparing the platform for production use
and completing the project. This involved configuring the
application for deployment, setting up the production environment,
and performing the final release. The goal was to make the
platform accessible to real users while ensuring stability
and maintainability.

The deployment process started with preparing the application
configuration for production. Environment variables were set
for production values, including database connections, caching
servers, and external service credentials. Debug mode was
disabled and error handling was configured to show user-friendly
messages rather than technical details.

The server infrastructure was configured to host the Laravel
application. This included setting up the web server, configuring
PHP runtime settings, and establishing the database connection.
The caching layer using Redis was verified to be functioning
correctly in the production environment.

The database migration system ensured that the production
database had the correct schema. Migration scripts that had
been developed throughout the project were executed against
the production database, creating all necessary tables and
indexes. Seed data was applied for foundational data like
categories and initial administrative accounts.

Deployment scripts automated the release process. These scripts
handled pulling the latest code, running migrations, clearing
and warming caches, and restarting background workers. This
automation reduced the risk of human error during deployment
and allowed for consistent reproducible releases.

Rollback procedures were established in case issues were
discovered after deployment. The version control system
maintained previous releases, allowing quick reversion if
necessary. Database migrations were designed to be reversible
where possible to facilitate rollbacks.

\begin{table}[H]
\begin{center}
\begin{tabular}{|m{3.5cm}|m{5.5cm}|}
\hline
\textbf{Deployment Task} & \textbf{Description} \\ \hline
Environment Configuration & Production settings for all services \\ \hline
Server Setup & Web server and PHP configuration \\ \hline
Database Preparation & Migrations and initial data setup \\ \hline
Release Automation & Automated deployment scripts \\ \hline
Rollback Procedures & Plan for reverting problematic releases \\ \hline
\end{tabular}
\caption{Deployment Components}
\label{tab:deployment_components}
\end{center}
\end{table}

Monitoring was configured to track the health and performance
of the production system. Log aggregation collected application
logs for analysis. Performance metrics were collected to track
response times and resource usage. Alert thresholds were set
to notify administrators of potential issues.

The user documentation was completed during this sprint. This
included API documentation for developers integrating with
the platform, user guides for the various user types, and
administrator documentation for managing the system.

A thorough final review verified that all planned features
were implemented and met the defined requirements. The product
backlog was reviewed to ensure all committed items were
complete. Any remaining items were moved to a future release
backlog for consideration in post-project maintenance.

Testing in the production environment verified that the
deployment was successful. End-to-end tests executed critical
user workflows to confirm everything worked as expected.
Load testing verified the system could handle expected
traffic levels.

The production launch was successful. Users could access
the platform and use all implemented features. The monitoring
system tracked system health and quickly identified any
initial issues that required attention.

\begin{figure}[H]
\begin{center}
\fbox{Add screenshot: Production dashboard / Deployment logs / Monitoring}
\end{center}
caption{Deployment Interface}
\label{fig:deployment_interface}
\end{figure}

% ============================================================
\section*{Conclusion and Perspectives}
\addcontentsline{toc}{section}{Conclusion and Perspectives}
\addcontentsline{toc}{subsection}{Conclusion}
% ============================================================

The PrixTunisix project was successfully completed with all
planned features implemented and deployed. The platform now
provides a comprehensive price comparison service that serves
multiple user types with their specific needs.

Visitors can browse products, search for items, and compare
prices across different merchants without creating accounts.
Authenticated clients gain access to personalized features
like favorites, price alerts, and wishlists that enhance their
shopping experience. Merchants can manage their product
offerings and track their performance on the platform.
Administrators have tools to manage users, validate products,
and monitor overall platform health.

The technical architecture proved effective throughout the
project. Laravel provided a solid foundation for rapid
development. The role-based authentication system scaled well
to support five distinct user types. The product catalog and
search functionality handled large datasets efficiently. The
scraping system automated data collection from merchant websites.

The SCRUM methodology worked well for managing the development
process. The eight sprints provided clear milestones and
accountability. Regular standups kept the team aligned. Sprint
reviews demonstrated progress to stakeholders. Retrospectives
identified areas for improvement.

The final product meets the requirements established at the
beginning of the project. Users can find the best prices on
products they want to purchase. The platform generates value
for both consumers seeking deals and merchants seeking customers.

% ============================================================
\subsection*{Future Work}
\addcontentsline{toc}{subsection}{Future Work}

While the current implementation provides a complete price
comparison platform, several enhancements could be made in
future iterations.

The mobile application could be developed to provide a native
experience for users who prefer accessing the platform through
their phones. This would require additional development beyond
the current responsive web interface.

Additional merchants could be integrated into the scraping
system, expanding the product catalog and increasing the
value provided to users. The matching system could be enhanced
with machine learning to improve accuracy.

Advanced features like price prediction could leverage
historical data to forecast future price movements. This
would help users decide whether to buy now or wait for a
better price.

Social features could enable users to share deals with friends
or collaborate on purchase decisions. This would increase
user engagement and grow the platform organically.

The integration with more external services could enhance
the platform's capabilities. Payment processing could enable
transactions directly on the platform rather than redirecting
to merchant sites.

These possibilities represent opportunities for continued
development beyond the initial project scope while building
upon the solid foundation that has been established.